%% clawmarc_wadl.tex — reviewer-facing prong for WADL 2026 (ACM sigconf).
%% Sole author: A. Kornai; AI assistance in Acknowledgments (ACM policy).
%% Parallel file: clawmarc.tex (clawxiv-ready prong, joint authorship).
%% The two are synced BY HAND; never refactor into shared \input bodies.
%% v1.rc2 (2026-09-07): restructured after the referee-report draft; DCMI classes; pilot.
\documentclass[sigconf]{acmart}

\setcopyright{none}
\settopmatter{printacmref=false}
\renewcommand\footnotetextcopyrightpermission[1]{}
\acmConference[WADL 2026]{Web Archiving and Digital Libraries Workshop}{October 16, 2026}{Dallas, TX, USA}
\acmYear{2026}
\copyrightyear{2026}

\usepackage{booktabs}

\begin{document}

\title{The 4096-byte catalog card: fixed-size, signed, content-bound description for archives that machines read}

\author{Andr\'as Kornai}
\orcid{0000-0001-6078-6840}
\email{kornai@sztaki.hu}
\affiliation{%
  \institution{SZTAKI Institute for Computer Science and Control}
  \city{Budapest}
  \country{Hungary}
}

\begin{abstract}
We describe clawmarc, a catalog card for digital artefacts that is exactly
4096 bytes long, binds by SHA-256 to the bytes it describes, carries a short
human-readable description alongside a machine-readable search payload,
records retrieval hints for several storage backends, and is signed by its
issuer with Ed25519. The card is a descriptor, not a proof of authenticity or
availability, and it is meant to live separate of the artefact; its artefact
classes are the twelve terms of the DCMI Type Vocabulary, and subject
description is left to the schemes a cataloger already uses. We give the
design and its rationale, place the card against MARC, Dublin Core, Linked
Data, WARC, Memento, and BagIt, and report what has been built: a public
specification with a C reference layout, a Python producer and inspector,
reference cards, a differential test in which two independent implementations
produced byte-identical cards on every pinned field, and a pilot that minted
cards for a small sample of a research archive and ranked them offline from
the cards alone.
\end{abstract}

\begin{CCSXML}
<ccs2012>
<concept>
<concept_id>10002951.10003317.10003318</concept_id>
<concept_desc>Information systems~Document representation</concept_desc>
<concept_significance>500</concept_significance>
</concept>
<concept>
<concept_id>10002951.10003317.10003338</concept_id>
<concept_desc>Information systems~Retrieval models and ranking</concept_desc>
<concept_significance>300</concept_significance>
</concept>
<concept>
<concept_id>10002951.10003227.10003392</concept_id>
<concept_desc>Information systems~Digital libraries and archives</concept_desc>
<concept_significance>500</concept_significance>
</concept>
</ccs2012>
\end{CCSXML}
\ccsdesc[500]{Information systems~Digital libraries and archives}
\ccsdesc[500]{Information systems~Document representation}
\ccsdesc[300]{Information systems~Retrieval models and ranking}

\keywords{catalog card, content addressing, metadata, fixed-size records, embeddings, DCMI Type Vocabulary, web archiving, interoperability}

\maketitle

\section{Introduction}

For about a century the unit of the library catalog was a card of
$7.5 \times 12.5$ centimetres. Its virtues were not typographic. A drawer of
cards could be counted, its space planned, a card pulled without disturbing its
neighbours, and a set of cards from one library merged into the drawers of
another; the Library of Congress sold printed cards precisely because every
recipient's drawers were the same size. The card did not need to be as long as
the book; it needed to be long enough to say what the book was, where it stood,
and who had described it. Krajewski's history of the format traces how much
intellectual infrastructure rested on that one fixed dimension
\cite{Krajewski:2011}.

MARC replaced the card with a variable-length record \cite{Avram:1975}, and
every later interchange format, from Dublin Core \cite{Kunze:2007} to OAI-PMH
\cite{Lagoze:2001} to the graph serializations of Linked Data
\cite{Kellogg:2020}, followed it in letting the record grow to fit the
description. This was the right choice for records that skilled catalogers
write and databases hold. It is less obviously right for the situation now
facing digital libraries and web archives: artefacts produced in volumes that
no cataloging staff can describe, a growing share of them produced by, or with,
language models; storage replicated across mirrors, institutional
repositories, and content-addressed peer-to-peer networks
\cite{Benet:2014,Alam:2016}; and description increasingly performed by machines
for machines, with the human reading a card only after a search has already
ranked it.

This paper describes a catalog card designed for that situation. A clawmarc
card is exactly 4096 bytes. It binds to the bytes of the artefact it describes
by SHA-256, so that a later copy found anywhere can be checked against it. It
carries a short human-readable description (title, abstract, keywords,
classification, and as much body prefix as fits) and, in the same card, a
machine-readable payload, in the current profile a 768-dimensional document
vector. It records retrieval hints for HTTP, IPFS, IPNS, and Swarm, and it is
signed by its issuer with Ed25519 \cite{Josefsson:2017}. It is a descriptor and
nothing more: it does not certify that the artefact is authentic, available, or
good, and it is meant to remain useful after the artefact has gone.

The format grew out of ClawXiv, an archival workflow for research produced
jointly by humans and language models \cite{Kornai:2026a}, where the need to
catalog machine-produced bundles at machine speed was felt
early. Section~\ref{sec:card} states the design.  Section~\ref{sec:rationale}
gives the reasons for its main choices, including the ones that run against
current practice. Section~\ref{sec:related} places the card among the formats
already in broad currency.  Section~\ref{sec:impl} reports what has been
built and measured, and Section~\ref{sec:governance} lays out the questions
that remain, most of them experimental.

\section{The card}
\label{sec:card}

\subsection{Layer model}

Three layers are kept apart. The \emph{card} is the 4096-byte descriptor. The
\emph{artefact} is the thing described. The \emph{catalog} is a collection of
cards together with indices, shards, heads, admission policy, and governance.
The format specifies only the first. A card can outlive its artefact;
authenticity and availability are resolved at the artefact and catalog layers.
This separation is what allows the card to be storage-neutral: cards and
artefacts may travel over local filesystems, HTTP mirrors, IPFS, Swarm, or an
institutional repository, and the card neither knows nor cares which.

\subsection{Layout}

The card is divided into a 1216-byte header of fixed-position fields, a
2816-byte \emph{arena} whose interpretation depends on the artefact class, and a
64-byte footer holding two CRC-32 checks. Table~\ref{tab:layout} summarizes the
header. Every multi-byte field lies on its natural alignment boundary; the
reference C structure compiles to exactly 4096 bytes with no padding and is
nevertheless declared packed, with compile-time assertions on the total size
and on the offsets of the fields most likely to drift, so that the wire layout
cannot silently change across compilers or across edits.

\begin{table}[t]
\caption{Header fields, in wire order (offsets in hex; the arena occupies
0x4c0--0xfbf, the footer 0xfc0--0xfff).}
\label{tab:layout}
\footnotesize
\begin{tabular}{@{}rrp{0.66\columnwidth}@{}}
\toprule
offset & size & field \\
\midrule
000 & 8   & magic \texttt{CXCC}, layout major/minor \\
008 & 8   & arena split, arena class, size class, flags \\
010 & 32  & card issued, work created, work revised, sequence \\
030 & 128 & schema, object, manifest, previous-card SHA-256 \\
0b0 & 96  & issuer public key, card signature (Ed25519) \\
110 & 132 & issuer card reference; access mode, encryption profile, key \\
194 & 88  & responsible ORCID; author, classification, URL counts; author digests \\
1ec & 8   & license identifier \\
1f4 & 272 & Swarm reference, IPFS CID, IPNS name, HTTP hint, locator-set digest \\
304 & 76  & embedding profile, four text lengths, text flags, text and embedding digests \\
350 & 32  & target card (bounded indirection) \\
370 & 336 & reserved (zero) \\
\bottomrule
\end{tabular}
\end{table}

\subsection{Bindings}

Four digests tie the card to the world. The \emph{object} digest is the SHA-256
of the artefact bytes: the card says ``I describe these bytes,'' and any later
copy can be checked. The \emph{schema} digest identifies the frozen
specification the producer followed, prose and reference header together, so a
consumer knows which layout it is reading. An optional \emph{manifest} digest
binds a build manifest, source listing, or migration record. The
\emph{previous-card} digest forms a supersession chain when a card is reissued.
None of these is a proof of authenticity: a valid signature proves only that a
named issuer made a signed statement about an object hash.

\subsection{The arena}

The arena is split at a producer-chosen offset into a machine payload and a
human text. For text artefacts the payload, when present, is a document vector
under a named \emph{embedding profile}. Three profiles are assigned:
profile~1 is the 384-dimensional English-language \texttt{bge-small-en-v1.5}
model \cite{Xiao:2023}, stored as little-endian binary16 in 768 bytes;
profiles~2 and~3 are the 768-dimensional multilingual
\texttt{bert-base-multilingual-cased} \cite{Devlin:2019} and
\texttt{paraphrase-multilin}\\\texttt{gual-mpnet-base-v2} \cite{Reimers:2020}, 1536
bytes in the same encoding, leaving 1280 bytes for text; profile~2 is the
default. Against profile~1 the trade is deliberate: half the human-readable
prefix buys a vector that serves the multilinguality of any larger
archive. The human
text holds title, abstract, keywords, classification, and whatever prefix of
the body fits; a digest over the full used text and a truncation flag make the
cut explicit. The classification field is free text and is where a producer
writes whatever subject scheme it uses: a Library of Congress class, a Dewey
number, an arXiv category, a schema.org type, or a JSON-LD fragment. For
images the payload may be a small visual reduction and the text a caption. For
opaque artefacts the split is zero and the text is a category descriptor; a
producer must not fabricate a document vector for bytes that contain no
running text.

\subsection{Classes}
\label{sec:classes}

The one-byte \emph{arena class} names what kind of thing is described. It is a
type, not a subject: the distinction an operating system draws between text,
image, executable, and archive, not the distinction a classification scheme
draws between endocrinology and epidemiology. Rather than mint such a
vocabulary, the format reuses one the profession already maintains.
Values 16--27 are the twelve classes of the DCMI Type Vocabulary
\cite{DCMI:2020}, assigned in the alphabetical order of that vocabulary
(Table~\ref{tab:classes}), so that the URI of a class is recoverable from its
number. Value 0 means no class asserted. Values 1 and 2 are structural and
belong to the format: an \emph{indirect} card describes an artefact that is
itself a concatenation of cards, readable at 4096-byte stride, and a
\emph{doubly indirect} card describes a concatenation of indirect cards.
Consumers descend at most two levels; this is aggregation with a bounded
depth, and it is how a catalog shards a million cards into a few hundred
artefacts without inventing a container format. Values 3--15 are reserved to
the format and 28--255 are unassigned.

\begin{table}[t]
\caption{Arena class values. DCMI terms carry the URI prefix
\texttt{http://purl.org/dc/dcmitype/}.}
\label{tab:classes}
\footnotesize
\begin{tabular}{@{}rl@{\hspace{1.5em}}rl@{}}
\toprule
0 & no class asserted        & 16 & Collection \\
1 & indirect (cards)         & 17 & Dataset \\
2 & doubly indirect          & 18 & Event \\
3--15 & reserved (structural) & 19 & Image \\
28--255 & unassigned         & 20 & InteractiveResource \\
      &                      & 21 & MovingImage \\
      &                      & 22 & PhysicalObject \\
      &                      & 23 & Service \\
      &                      & 24 & Software \\
      &                      & 25 & Sound \\
      &                      & 26 & StillImage \\
      &                      & 27 & Text \\
\bottomrule
\end{tabular}
\end{table}

\section{Design rationale}
\label{sec:rationale}

\subsection{Fixed size}

Nothing in current practice argues for a fixed record size, and the choice
needs defending. The defence has three parts, none of them about saving bytes
or parser cycles; a modern machine parses variable-length records fast enough.

First, the operations a catalog performs on cards, as distinct from the
operations it performs on artefacts, are exactly the operations that fixed
size makes trivial. A shard of cards is addressed by ordinal; it is
memory-mapped and scanned at constant stride; a card is validated by a single
length check before any field is parsed; a corrupted card damages one slot and
not its neighbours. Ranking a corpus of cards by a query vector is one
matrix product over a dense array that the shard already is. Each of these has
an equivalent for a JSON record set, and each equivalent needs an index, a
parser, or a pointer chain whose failure is a failure of the whole.

Second, planning. An archive that knows its holdings knows its catalog's size
to the byte: a million artefacts is four gigabytes of cards, whatever the
artefacts are. Replication, sharding, and the decision to keep a full mirror
of another institution's cards are decisions about a known quantity.

Third, the card is not the record. A variable-length description belongs in
the artefact layer, or in a manifest the card binds by digest, and both may be
as long as they like. The card is the thing that must be everywhere, quickly,
in a form that can be trusted without being trusted much. The 3-by-5 card said
less than the title page; the point of it was that it said enough and fitted
the drawer.

Why 4096? It is one common filesystem block and one common virtual-memory
page, so a card is read or written in one unit and never straddles two. It
holds the digests, keys, signature, and locators the header needs, a
768-dimensional vector in half precision, and over a kilobyte of
prose. Alignment within the card is the same argument at a smaller scale, and
the compile-time assertions make it cheap to keep. The drawbacks are trivial:
text, such as extremely verbose abstracts, may get truncated; short artefacts
waste most of the arena (but note that ext4, NTFS, APFS, and XFS all default
to 4k frag size). 

\subsection{A machine payload beside the human text}

To defend against false positives, cataloging work increasingly aims at high
precision, often by combining orthogonal descriptors. Suppose our interest is
with endometriosis in Peru. A paper about uterine diseases among the Quechua
could be highly relevant, but will not surface based on a keyword match.  The
machine payload adds a document vector that carries such loose associations,
broader and narrower alike, and it does so for queries not anticipated when
the card was written. It is stored on the card so that a catalog can rank
without fetching the artefact and without a central index: a drawer of cards
and a query vector is a search engine.

The proposed card does not help with quality issues. Tens of thousands of
papers are genuinely about quantum gravity; they are not false positives, and
they vary in quality by orders of magnitude. Trade coverage of any industry is
vast, on topic, and shot through with interest. What the card can do is make
the description honest about its own provenance: the issuer is named and has
signed; the embedding profile is named; the text that was embedded is
digested; and the extraction and reduction policy is recorded in a manifest
the card binds. Two producers describing the same bytes will differ in their
text and therefore in their vectors, and consumers compare vectors by cosine
similarity rather than byte identity. Section~\ref{sec:impl} reports how far
that variation reaches in practice.

\subsection{Embedding profiles and their obsolescence}
\label{sec:profiles}

Embedding models age quickly, and a format that froze one would age with it.
The card does not. The profile field is sixteen bits and names model,
dimension, and encoding together; the arena split is per card, so moving from
a 384-dimensional to a 768-dimensional profile costs only text, a
1024-dimensional half-precision vector fits the same way, and a binary or
8-bit quantization fits with room to spare.
When a catalog moves to a new profile it reissues cards along the supersession
chain, and a consumer holding cards under two profiles ranks each population
separately, as it would two indices. What the format fixes is where the vector
lives and how it is described, not which model made it.

Which model made it is nonetheless a question the catalog must be able to
answer exactly, and a colloquial model name does not: the same weights
circulate in several serializations, from several hubs, no two files
byte-identical, and the vector a consumer recomputes depends further on the
tokenizer files. The proposal is that the model be incorporated into the
catalog itself. The weights and tokenizer of a profile's model, in one
canonical serialization, are an artefact like any other; each profile
references a clawmarc card minted for that artefact, which binds the exact
bytes by digest, records the reduction in its manifest, carries mirror
locators, and is signed by the profile's issuer. A catalog that archives the
model beside the cards it validates no longer depends on any distribution
hub's longevity, and a consumer asking which model profile~2 names gets an
answer with the same shape, and subject to the same checks, as any other
artefact in the drawer.

\subsection{Content binding and format migration}

A digital library normalizes formats: a PDF is migrated to PDF/A, a TIFF is
rewrapped, headers change. Every such act changes the bytes, and a card bound
to the old bytes no longer matches the new file. This is the intended
behaviour of a fixity binding, and it is exactly the property that lets a
reader notice that an archive has changed under them \cite{Lerner:2017,
Aturban:2019}; but it means the workflow must be stated. A migrated artefact
is a new artefact. It gets a new card, whose previous-card digest points at the
old one and whose manifest digest binds the migration record, so that the
chain carries the history and the old card remains true of the bytes it
described. This is the same discipline BagIt applies to a bag's payload
manifest \cite{Kunze:2018}, lifted from the package to the catalog.

\subsection{Publication posture}

A card is a publication act. Even when the artefact is withheld, a card may
disclose enough to identify, summarize, retrieve, classify, or compare it; a
document vector reveals gist-level information about the text it represents
and is not anonymization. The producer therefore runs under an explicit
content policy, shipped as an editable JSON profile, that says whether a text
prefix, a vector, a thumbnail, OCR, or an archive member listing may be stored,
and the policy is recorded in the card's manifest. A profile suited to a public
science library allows all of these; a profile for a directory of mixed-rights
material allows almost none and records hashes only. The same institutional,
legal, and ethical review that applies to a finding aid applies to a card, and
the format is built so that a producer can obey it.

\section{Relation to existing practice}
\label{sec:related}

Dublin Core \cite{Kunze:2007} is the closest relative: a deliberately small
element set, meant for machine harvesting and cross-domain use, whose Type
Vocabulary the card adopts outright. The card differs in being fixed-size,
signed, and bound to the bytes described; it can be read as a signed,
content-bound carrier for a Dublin-Core-sized description plus a vector.
OAI-PMH \cite{Lagoze:2001} moves such records between repositories on request;
a shard of cards is the same idea with the transport removed, so that a mirror
is a copy of a file.

The metadata community's direction of travel for twenty years has been toward
graphs: Linked Data, RDF, JSON-LD \cite{Kellogg:2020}, and the shared
vocabulary of schema.org \cite{Guha:2016}. The card is not a rival to that
work and would lose if it were. A graph is extensible, mergeable, and
expressive; a card is none of these. What a card offers a graph is an
envelope: a fixed-size, signed, content-bound unit that can carry a JSON-LD
fragment or a schema.org type in its classification text, can be replicated
without a triple store, and can be ranked without being parsed. Signposting
\cite{VanDeSompel:2015} makes a related bargain for the live web, exposing a
few typed links so that a machine can navigate a scholarly object without
understanding it; a card is a signpost that survives offline.

In web archiving, WARC \cite{ISO:2017} is the record of capture, and the card
is not a substitute for it; a card can describe a WARC file, or a member of
one, as it can describe anything with bytes. InterPlanetary Wayback
\cite{Alam:2016} put WARC payloads on IPFS and is the nearest precedent for
content-addressed replication of archival holdings; a card carries the IPFS
CID and Swarm reference of what it describes and can serve as a signed,
portable index entry for such a store. Memento \cite{VanDeSompel:2013} answers
``what did this URI look like at time~$t$'' from a live archive; a card
answers the narrower question ``which bytes did this issuer describe, and where
might they be,'' and does so offline. The fixity work of Aturban, Alam,
Nelson, and Weigle \cite{Aturban:2019}, and the demonstration by Lerner, Kohno,
and Roesner that the archived web can be altered from the present
\cite{Lerner:2017}, motivate the object binding directly: a signed digest
issued at cataloging time and held outside the archive is what lets a later
reader notice the change. BagIt \cite{Kunze:2018} does this for a package;
the card does it for a catalog entry.

\section{Implementation and experience}
\label{sec:impl}

The public bundle contains the specification as an Internet-Draft in Markdown,
XML, text, and HTML; a C header that is the normative layout; a Python
producer, inspector, and general-file cataloger; three editable content
policies; batch scripts that mint indirect and doubly indirect collection
cards over a directory tree; and two reference cards with their JSON
metadata.\footnote{\url{https://github.com/kornai/clawmarc}}

\paragraph{Reference layout.} The header declares the structure packed and
asserts at compile time that it is 4096 bytes, that the arena begins at 0x4c0,
that the footer begins at 0xfc0, and that every multi-byte field sits on its
natural boundary. A producer in any language can be checked against it by
compiling one file.

\paragraph{Producer and inspector.} The Python producer detects type by
content, extracts text with version-pinned tools (pdftotext, pandoc, detex),
normalizes to Unicode NFC, computes the document vector when policy allows,
packs the card, signs it, and emits a JSON sidecar recording the extraction
and reduction policy and the exact text that was embedded. The inspector
validates length, magic, split, reserved-zero fields, both CRCs, and the
signature, and recomputes the object binding when the artefact is at hand. A
deterministic issuer seed lets two implementations following the same policy
byte-compare their output.

\paragraph{Differential test.} Two implementations were written independently
by two different language-model systems working from the same specification,
without sight of each other's code, and each minted a card for the same
artefact. The cards agreed byte for byte on every field the specification
pins: object and schema digests, class, split, embedding profile, ORCID,
license, chain and indirection fields, and title. They disagreed on every field
the specification had left to producer judgment: which of several valid mirror
locators to record, whether to populate work timestamps from the bundle, and
what to put in the abstract and keywords, where one producer wrote a gist and
the other scraped the source. The embeddings told the same story in numbers.
Two runs of one implementation agreed to cosine 0.9999999999. The same text
segments run through the other implementation's pipeline agreed to 0.989,
which measures how much the construction of the embedding input was unpinned.
The two cards as minted agreed to 0.932, which measures how much the text
differed. The specification was tightened on the first two points; the third
is the residual that consumers must expect, and the reference cards are
published with that expectation stated.

\paragraph{Pilot.} Twenty PDFs were drawn at
random from the 3,295 in the author's paper archive and minted on a laptop with text
extraction, a document vector under profile~2, and a body prefix under the
public-library policy. All twenty produced valid cards of class Text.
Extraction took 0.08~s per file on average and minting 5.9~s, most of it
interpreter start and model load that a batch producer pays once. The embedded
text averaged 74,000 characters and is the entire extracted body: the vector
is computed over the full text, however long, while the card's human-readable
prefix is truncated at its 1280 bytes, as the design expects. Over the
resulting 80~KB of cards, and nothing else, retrieval behaved as intended.
Fillmore's 1968 case grammar paired with Culicover and Jackendoff's
\emph{Simpler Syntax} (0.983); a 2015 study of recurrent network
architectures, a 2017 paper on neural architecture search, and a 2023 paper on
quantizing large language models formed a cluster with pairwise cosine at
least 0.977; and the document unlike anything else in the sample, a 1958
philosophy paper, had no neighbour above 0.94. Pairwise cosine over all 190
pairs ranged from 0.84 to 0.99 with mean 0.95; how much of that compression
is the model rather than the sample is measured next. A keyword scan of the
cards' text regions for \emph{grammar} returned the three linguistics papers
and nothing else.

\paragraph{Model comparison.} Four models were run over the same twenty
documents with the pipeline otherwise fixed: full extracted text, chunks at
the model's native window, token-count weighted mean, final L2
normalization. Throughput on the laptop was flat across the three
768-dimensional multilingual candidates, 10,900--11,700 tokens per second,
and 32,700 for the 384-dimensional profile~1. Everything else differed. The
untuned \texttt{xlm-roberta-base} collapsed: all 190 pairwise cosines fell
between 0.996 and 1.000, the entire ranking signal squeezed into the third
decimal place, and no profile is assigned to it. Multilingual BERT
(profile~2) spread the same pairs over 0.839--0.991; the English-only
sentence-tuned profile~1 model over 0.592--0.959; and the multilingual
sentence-tuned profile~3 model over 0.156--0.903, with the sample's two
singleton documents --- the 1958 philosophy paper and a paper on
web-archive security --- isolated below 0.57 while cluster neighbours sit
near 0.9. Every usable model finds the same linguistics and
neural-computation clusters and the same isolates, but the agreement has
structure: the two sentence-tuned models, which share no architecture,
training data, or language coverage, rank the 190 pairs in the closest
accord of any pair of models (Spearman 0.90), while the two untuned
encoders agree that well only with each other (0.88), and tuned--untuned
agreement runs at 0.80 and below. The natural reading is that the tuned
models approximate a common similarity which the untuned ones see through
noise --- and with very different margins, which is what a consumer ranking
vectors recomputed in unknown runtimes will spend. Quantization, by
contrast, costs nothing measurable: rounding the vectors to binary16, the
card's storage encoding, left rank order identical for all four models
(Spearman 1.000, every nearest neighbour unchanged), the per-component
rounding error averaging out across the hundreds of dimensions of the dot
product.

\section{Future work}
\label{sec:governance}

The artefact classes now reuse a vocabulary the Dublin Core Metadata
Initiative maintains, and the structural values are fixed by the
specification, so the class byte needs no registrar. Two enumerations still
do: the embedding profiles, which will grow as models change, and the
encryption profiles. The specification leaves most of each space unassigned,
names no authority, and asks producers not to consume unassigned values in
public cards except by agreement among the catalogs that will read them.
Three resolutions are open: a registry lodged with an existing body, profiles
named by URI in a manifest and resolved rather than enumerated, or continued
enumeration by agreement while producers remain few. The format has been
written so that any of the three can be adopted later without a layout
change, and the choice will be made by the number of producers rather than by
argument.

\paragraph{Model identity.} The model cards of
Section~\ref{sec:profiles} make the identification of a profile's model
exact, but one quantity in them is recorded rather than pinned: the runtime.
A recomputed vector depends on the pooling code and the floating-point
environment as well as on the weights, and while the differential test above
bounds reproducibility within one environment at ten decimal places, drift
across environments is unmeasured. Measuring it is the next experiment, run
the way the differential test was run: the same model card, several
environments, a cosine agreement matrix. If cross-environment drift lands
near the within-environment bound, digest-pinning the weights suffices and
the runtime stays advisory; if not, the profile must also name the runtime,
and the experiment will have said so before a reviewer did.

\paragraph{Choice of model.} Which model deserves the next profile is an
empirical question with three axes: retrieval quality, language coverage, and
the durability of the weight distribution, which is not a hypothetical: even
the most widely used checkpoints are today served by a small number of hubs,
rate-limited, mirrored unevenly, and republished in mutually incompatible
serializations. The comparison in Section~\ref{sec:impl} prices the quality
axis on one small English sample: multilingual BERT, the most widely
replicated checkpoint in existence, keeps a workable ranking margin; the
sentence-tuned model has several times that margin but one canonical
distribution point; and the large long-context embedders, not yet measured,
cost dimensions the arena must pay for in text. The model cards of
Section~\ref{sec:profiles} remove the distribution axis entirely; the
remaining two are settled by running the harness at scale, on multilingual
samples with relevance judgments, not by taste.

\paragraph{Reduction.} The reduction from an unbounded text to one vector is
recorded in the manifest but left to producer choice: chunk size, chunk
weighting, and pooling are all defensible in several ways, and the 0.989
agreement figure in the differential test measures exactly the cost of
leaving them unpinned. Candidate reductions are compared by the same offline
IR harness, and the specification will pin one when the comparison says
something.

\section{Conclusion}

The format described here assumes that the descriptor of a digital artefact
should have the same shape everywhere, and derives the rest: content binding
so the descriptor can be checked, a signature so it can be attributed,
locators so it can be followed, a machine payload so it can be ranked without
being fetched, and a fixed size so that all of this can be replicated and
scanned as a dense array. The specification, reference implementation, and
reference cards are public; the measurements reported are small and are
reported as such; and the questions left open are stated with either the
experiment or the community that will decide them.

\begin{acks}
We thank Eliot Lear (Cisco) for penetrating criticism of an earlier draft.
The specification and the reference implementation were developed with AI
assistance from GPT-5 Codex (OpenAI) and Claude Opus (Anthropic), which also
performed the two independent implementations of the differential test and
adversarial review of each other's work; Gemini 3.1 Pro (Google) ran a first
pass of the pilot that exposed a missing PDF branch in the cataloger. Claude
Fable 5.1 (Anthropic) drafted this paper from the specification, the project
record, and the author's correspondence, under the author's direction; the
author is responsible for its content. This disclosure follows the ACM policy
on authorship and generative AI.
\end{acks}

\bibliographystyle{ACM-Reference-Format}
\bibliography{clawmarc}

\end{document}
